\documentclass[10pt]{article}
\usepackage[utf8]{inputenc}
\usepackage[T1]{fontenc}
\usepackage{hyperref}
\hypersetup{colorlinks=true, linkcolor=blue, filecolor=magenta, urlcolor=cyan,}
\urlstyle{same}
\usepackage{amsmath}
\usepackage{amsfonts}
\usepackage{amssymb}
\usepackage[version=4]{mhchem}
\usepackage{stmaryrd}
\usepackage{caption}

\title{Data Compression Course Project Implementation of BZip2 Compression Algorithm }

\author{Course Instructor: Dr. Faisal Aslam}
\date{}


\DeclareUnicodeCharacter{2192}{\ifmmode\rightarrow\else{$\rightarrow$}\fi}

\begin{document}
\maketitle
\captionsetup{singlelinecheck=false}
Date: April 14, 2026

\section*{Contents}
1 Introduction ..... 3\\
1.1 Project Overview ..... 3\\
1.2 Project Timeline ..... 3\\
1.3 Evaluation Criteria ..... 3\\
2 System Architecture ..... 4\\
2.1 Overall Block Diagram ..... 4\\
2.2 Configuration Management ..... 4\\
3 Stage 1 Implementation (1.5 Weeks) ..... 5\\
3.1 Block Division and File Handling ..... 5\\
3.1.1 Requirements ..... 5\\
3.1.2 Data Structures ..... 5\\
3.1.3 Function Prototypes ..... 5\\
3.2 RLE-1 Implementation ..... 6\\
3.2.1 Description ..... 6\\
3.2.2 Function Prototypes ..... 6\\
3.3 Burrows-Wheeler Transform (BWT) ..... 6\\
3.3.1 Description ..... 6\\
3.3.2 Data Structures ..... 7\\
3.3.3 Function Prototypes ..... 7\\
3.4 Stage 1 Evaluation ..... 7\\
4 Stage 2 Implementation (6 Days) ..... 8\\
4.1 Move-to-Front (MTF) Transform ..... 8\\
4.1.1 Description ..... 8\\
4.1.2 Function Prototypes ..... 8\\
4.2 RLE-2 Implementation ..... 8\\
4.2.1 Description ..... 8\\
4.2.2 Function Prototypes ..... 8\\
4.3 Stage 2 Evaluation ..... 9\\
5 Stage 3 Implementation (1 Week) ..... 9\\
5.1 Canonical Huffman Coding ..... 9\\
5.1.1 Description ..... 9\\
5.1.2 Data Structures ..... 9\\
5.1.3 Function Prototypes ..... 9\\
6 Build System and Platform Support ..... 11\\
6.1 Makefile Requirements ..... 11\\
7 Performance Evaluation and Benchmarking ..... 11\\
7.1 Benchmark Datasets ..... 11\\
7.2 Performance Metrics ..... 11\\
7.3 Automated Results Generation ..... 12\\
8 Extra Features and Enhancements (10\% Marks) ..... 12\\
8.1 Enhanced RLE Versions (Extra Marks) ..... 12\\
8.2 Suffix Array-based BWT (Extra Marks) ..... 12\\
8.3 Alternative Entropy Coding ..... 12\\
9 Submission Guidelines ..... 13\\
9.1 GitHub Repository Structure ..... 13\\
9.2 \href{http://README.md}{README.md} Requirements ..... 13\\
9.3 Graph Generation ..... 13\\
10 Important Notes for Students ..... 14\\
11 Conclusion ..... 14

\section*{1 Introduction}
\subsection*{1.1 Project Overview}
This project involves implementing a simplified version of the BZip2 compression algorithm. The project spans three weeks and is completed by a team of three individuals. The implementation includes all major components of BZip2: Run-Length Encoding (RLE), Burrows-Wheeler Transform (BWT), Move-to-Front (MTF) transform, and Huffman coding.

Important: Students must write their own code. Only function prototypes and struct definitions are provided as guidance. No implementation code is given.

\subsection*{1.2 Project Timeline}
\begin{itemize}
  \item Total Duration: 3 Weeks
  \item Stage 1: 1.5 Weeks (50\% of project grade)
  \item Stage 2: 6 Days ( $25 \%$ of project grade)
  \item Stage 3: 1 Week ( $25 \%$ of project grade)
\end{itemize}

\subsection*{1.3 Evaluation Criteria}
\begin{itemize}
  \item Each stage includes evaluation of both encoding and decoding
  \item Extra features: $10 \%$ additional marks
  \item Performance beating latest BZip2: $10 \%$ additional marks (max 2 groups)
\end{itemize}

\section*{2 System Architecture}
\subsection*{2.1 Overall Block Diagram}
\begin{verbatim}
Encoding Pipeline:
================
Input File
    |
    v
Block Division
    |
    v
RLE-1 Encoding
    |
    v
BWT
    |
    v
MTF
    I
    v
RLE-2 Encoding
    |
    v
Huffman Encoding
    I
    v
Compressed Output
\end{verbatim}

Figure 1: Encoding Pipeline Architecture

\subsection*{2.2 Configuration Management}
The system uses a configuration file named config.ini with the following structure:

\begin{verbatim}
[General]
block_size = 500000 # Bytes ( 100KB to 900KB)
rle1_enabled = true
bwt_type = matrix # matrix or suffix_array
mtf_enabled = true
rle2_enabled = true
huffman_enabled = true
[Performance]
benchmark_mode = false
output_metrics = true
[Paths]
input_directory = ./benchmarks/
\end{verbatim}

\begin{verbatim}
output_directory = ./results/
\end{verbatim}

Listing 1: Configuration File Format

\section*{3 Stage 1 Implementation (1.5 Weeks)}
\subsection*{3.1 Block Division and File Handling}
\subsection*{3.1.1 Requirements}
\begin{itemize}
  \item Support for very large files
  \item Configurable block size $(100 \mathrm{~KB}$ - 900 KB$)$
  \item Configuration file integration
\end{itemize}

\subsection*{3.1.2 Data Structures}
\begin{verbatim}
/*
    * Structure to hold a single block of data
    */
typedef struct {
        unsigned char *data; // Pointer to block data
        size_t size; // Current size of block
        size_t original_size; // Original size before compression
} Block;
/*
    * Structure to manage multiple blocks
    */
typedef struct {
        Block *blocks; // Array of blocks
        int num_blocks; // Number of blocks
        size_t block_size; // Configurable block size
} BlockManager;
\end{verbatim}

Listing 2: Block Division Structures (Provided)

\subsection*{3.1.3 Function Prototypes}
\begin{verbatim}
/*
    * Reads input file and divides into blocks
    * @param filename: Input file path
    * @param block_size: Size of each block in bytes
    * @return: BlockManager structure containing all blocks
    */
BlockManager* divide_into_blocks(const char *filename, size_t
        block_size);
/*
    * Reassembles blocks back into original file
    * Oparam manager: BlockManager containing processed blocks
    * @param output_filename: Path for output file
    * @return: 0 on success, -1 on failure
\end{verbatim}

\begin{verbatim}
*/
int reassemble_blocks(BlockManager *manager, const char *
    output_filename) ;
/*
* Frees memory allocated for BlockManager
* Cparam manager: Pointer to BlockManager to free
*/
void free_block_manager(BlockManager *manager);
\end{verbatim}

Listing 3: Block Management Functions

\subsection*{3.2 RLE-1 Implementation}
\subsection*{3.2.1 Description}
Run-Length Encoding replaces sequences of repeated characters with a count and the character.

\section*{Example:}
Input: $A B B C C C C D$\\
Output: A3B4C1D

\subsection*{3.2.2 Function Prototypes}
\begin{verbatim}
/*
    * Encodes data using Run-Length Encoding
    * @param input: Input byte array
    * @param len: Length of input array
    * @param output: Output buffer (must be pre-allocated)
    * @param out_len: Pointer to store output length
    */
void rle1_encode(unsigned char *input, size_t len,
        unsigned char *output, size_t *out_len);
/*
    * Decodes RLE-1 encoded data
    * @param input: Encoded byte array
    * @param len: Length of encoded data
    * @param output: Output buffer for decoded data
    * @param out_len: Pointer to store decoded length
    */
void rle1_decode(unsigned char *input, size_t len,
        unsigned char *output, size_t *out_len);
\end{verbatim}

Listing 4: RLE-1 Functions

\subsection*{3.3 Burrows-Wheeler Transform (BWT)}
\subsection*{3.3.1 Description}
The matrix-based BWT creates all cyclic rotations of the input string, sorts them lexicographically, and outputs the last column.

\subsection*{3.3.2 Data Structures}
\begin{verbatim}
/*
    * Structure for rotation in BWT
    */
typedef struct {
        char *rotation; // Rotation string
        int index; // Original index
} Rotation;
\end{verbatim}

Listing 5: BWT Structures

\subsection*{3.3.3 Function Prototypes}
\begin{verbatim}
/*
    * Compares two rotations for sorting
    * @param a: First rotation
    * @param b: Second rotation
    * @return: Comparison result (-1, 0, 1)
    */
int compare_rotations(const void *a, const void *b);
/*
    * Forward BWT transform
    * @param input: Input byte array
    * @param len: Length of input
    * @param output: Output buffer for BWT result
    * @param primary_index: Pointer to store primary index
    */
void bwt_encode(unsigned char *input, size_t len,
        unsigned char *output, int *primary_index);
/*
    * Inverse BWT transform
    * Cparam input: BWT encoded data
    * @param len: Length of encoded data
    * @param primary_index: Primary index from encoding
    * @param output: Output buffer for original data
    */
void bwt_decode(unsigned char *input, size_t len,
        int primary_index, unsigned char *output);
\end{verbatim}

Listing 6: BWT Functions

\subsection*{3.4 Stage 1 Evaluation}
Both encoding and decoding for Stage 1 components (RLE-1 and BWT) must be fully functional. The evaluation tests:

\begin{itemize}
  \item Correctness of RLE-1 encoding/decoding
  \item Correctness of BWT forward/inverse transform
  \item Block division and reassembly
  \item Configuration file parsing
\end{itemize}

\section*{4 Stage 2 Implementation (6 Days)}
\subsection*{4.1 Move-to-Front (MTF) Transform}
\subsection*{4.1.1 Description}
MTF encodes data by maintaining a list of symbols and outputting the index of each symbol, then moving it to the front.

\subsection*{4.1.2 Function Prototypes}
\begin{verbatim}
/*
    * Forward MTF transform
    * Cparam input: Input byte array
    * @param len: Length of input
    * @param output: Output buffer for MTF indices
    */
void mtf_encode(unsigned char *input, size_t len,
        unsigned char *output);
/*
    * Inverse MTF transform
    * @param input: MTF encoded indices
    * @param len: Length of input
    * @param output: Output buffer for decoded data
    */
void mtf_decode(unsigned char *input, size_t len,
        unsigned char *output);
\end{verbatim}

Listing 7: MTF Functions

\subsection*{4.2 RLE-2 Implementation}
\subsection*{4.2.1 Description}
RLE-2 specifically targets the output of MTF, which contains many zeros and small numbers. This is a specialized RLE for the MTF output.

\subsection*{4.2.2 Function Prototypes}
\begin{verbatim}
/*
    * Encodes MTF output using specialized RLE
    * @param input: MTF output array
    * @param len: Length of input
    * @param output: Output buffer
    * @param out_len: Pointer to store output length
    */
void rle2_encode(unsigned char *input, size_t len,
        unsigned char *output, size_t *out_len);
/*
    * Decodes RLE-2 encoded data
    * Oparam input: RLE-2 encoded data
    * @param len: Length of encoded data
\end{verbatim}

\begin{verbatim}
* Oparam output: Output buffer for MTF data
* @param out_len: Pointer to store output length
*/
void rle2_decode(unsigned char *input, size_t len,
    unsigned char *output, size_t *out_len);
\end{verbatim}

Listing 8: RLE-2 Functions

\subsection*{4.3 Stage 2 Evaluation}
The complete pipeline (RLE-1 → BWT → MTF → RLE-2 ) must work for both encoding and decoding.

\section*{5 Stage 3 Implementation (1 Week)}
\subsection*{5.1 Canonical Huffman Coding}
\subsection*{5.1.1 Description}
Canonical Huffman coding ensures efficient tree representation by storing only code lengths, not the full tree.

\subsection*{5.1.2 Data Structures}
\begin{verbatim}
/*
    * Huffman code representation
    */
typedef struct {
        unsigned short code; // Huffman code
        unsigned char length; // Code length in bits
} HuffmanCode;
/*
    * Huffman tree node
    */
typedef struct Node {
        unsigned char symbol; // Byte value (0-255)
        int freq; // Frequency count
        struct Node *left; // Left child
        struct Node *right; // Right child
} HuffmanNode;
\end{verbatim}

Listing 9: Huffman Data Structures

\subsection*{5.1.3 Function Prototypes}
\begin{verbatim}
/*
    * Builds Huffman tree from frequency counts
    * @param frequencies: Array of 256 frequency counts
    * @param root: Pointer to store root of Huffman tree
    */
void build_huffman_tree(int *frequencies, HuffmanNode **root);
\end{verbatim}

\begin{verbatim}
/*
    * Generates canonical Huffman codes from tree
    * @param root: Root of Huffman tree
    * @param codes: Array of 256 to store generated codes
    */
void generate_canonical_codes(HuffmanNode *root,
                HuffmanCode *codes);
/*
    * Encodes data using Huffman coding
    * Oparam input: Input byte array
    * Cparam len: Length of input
    * Cparam output: Output buffer for compressed data
    * Cparam out_len: Pointer to store output length
    */
void huffman_encode(unsigned char *input, size_t len,
                unsigned char *output, size_t *out_len);
/*
    * Decodes Huffman encoded data
    * @param input: Huffman encoded data
    * @param len: Length of encoded data
    * @param output: Output buffer for decoded data
    * @param out_len: Pointer to store output length
    */
void huffman_decode(unsigned char *input, size_t len,
                unsigned char *output, size_t *out_len);
/*
    * Writes Huffman header (code lengths) to output
    * @param codes: Array of Huffman codes
    * @param output: Output buffer
    * @param out_len: Pointer to update with header size
    */
void write_header(HuffmanCode *codes, unsigned char *output,
            size_t *out_len);
/*
    * Encodes data using generated codes
    * @param input: Input data
    * @param len: Length of input
    * @param codes: Huffman codes for each symbol
    * Oparam output: Output buffer
    * @param out_len: Pointer to update with encoded size
    */
void encode_data(unsigned char *input, size_t len,
            HuffmanCode *codes, unsigned char *output,
            size_t *out_len);
\end{verbatim}

Listing 10: Huffman Functions

\section*{6 Build System and Platform Support}
\subsection*{6.1 Makefile Requirements}
The project must include a cross-platform Makefile supporting both Linux and Windows. The Makefile should have the following targets:

\begin{verbatim}
# Required Makefile targets:
# all - Compile the complete project
# clean - Remove object files and executable
# windows - Cross-compile for Windows platform
# Example structure (students must complete):
CC = gcc
CFLAGS = -Wall-02
TARGET = bzip2_impl
SOURCES = main.c rle.c bwt.c mtf.c huffman.c config.c
# Platform detection and flags must be implemented
# Windows cross-compilation must be supported
\end{verbatim}

Listing 11: Makefile Structure (Students must implement)

\section*{7 Performance Evaluation and Benchmarking}
\subsection*{7.1 Benchmark Datasets}
Students must use the following publicly available benchmark files for testing:

\begin{table}[h]
\begin{center}
\captionsetup{labelformat=empty}
\caption{Table 1: Benchmark Datasets}
\begin{tabular}{|l|l|l|}
\hline
Dataset & Size & Type \\
\hline
Canterbury Corpus & Various & Mixed \\
\hline
Calgary Corpus & Various & Mixed \\
\hline
Silesia Corpus & Various & Mixed \\
\hline
Large Text Files & $10 \mathrm{MB}-100 \mathrm{MB}$ & Text \\
\hline
Binary Files & $10 \mathrm{MB}-100 \mathrm{MB}$ & Binary \\
\hline
\end{tabular}
\end{center}
\end{table}

\subsection*{7.2 Performance Metrics}
Performance is evaluated using a weighted formula:

$$
\text { Score }=w_{1} \cdot \frac{C_{\mathrm{ref}}}{C}+w_{2} \cdot \frac{S}{S_{\mathrm{ref}}}
$$

Where:

\begin{itemize}
  \item $C=$ Compression ratio achieved
  \item $C_{\text {ref }}=$ Reference compression ratio (bzip2)
  \item $S=$ Speed (MB/s)
  \item $S_{\text {ref }}=$ Reference speed (bzip2)
  \item $w_{1}, w_{2}=$ Weights (to be determined by instructor)
\end{itemize}

\subsection*{7.3 Automated Results Generation}
Students must implement a script to automatically generate results in results.csv with the following format:

\begin{verbatim}
# results.csv must contain:
File,Size,BlockSize,CompressionRatio,Time,Memory
# Students must implement their own benchmarking script
\end{verbatim}

Listing 12: Required Results Format

\section*{8 Extra Features and Enhancements (10\% Marks)}
These will be $10 \%$ marks of the overall project grade.

\subsection*{8.1 Enhanced RLE Versions (Extra Marks)}
Students can implement:

\begin{itemize}
  \item RLE with run-length thresholds: Only encode runs longer than a configurable threshold
  \item Adaptive RLE: Dynamically adjust encoding based on data characteristics
  \item RLE with entropy coding: Combine with simple entropy coding for better compression
\end{itemize}

\subsection*{8.2 Suffix Array-based BWT (Extra Marks)}
Faster implementation using suffix array algorithm:

\begin{verbatim}
/*
    * Builds suffix array for efficient BWT
    * @param text: Input text
    * @param n: Length of text
    * Creturn: Suffix array
    */
int *build_suffix_array(unsigned char *text, int n);
\end{verbatim}

Listing 13: Suffix Array BWT Prototype

\subsection*{8.3 Alternative Entropy Coding}
\begin{itemize}
  \item ANS (Asymmetric Numeral Systems): Modern entropy coder with nearoptimal compression
  \item Range Coding: Arithmetic coding alternative with integer precision
\end{itemize}

\section*{9 Submission Guidelines}
\subsection*{9.1 GitHub Repository Structure}
\begin{verbatim}
project-bzip2/
|-- src/
| |-- main.c
| |-- rle.c
| |-- bwt.c
| |-- mtf.c
| |-- huffman.c
| '-- config.c
|-- include/
| '-- *.h
|-- benchmarks/
| '-- (test files)
|-- results/
| '-- (output files)
|-- Makefile
|-- config.ini
|-- README.md
'-- report.pdf
\end{verbatim}

\subsection*{9.2 README.md Requirements}
The README must include:

\begin{itemize}
  \item Complete feature description
  \item Implementation details for each stage (written by students)
  \item Build instructions for Linux and Windows
  \item Usage examples
  \item Performance results and graphs
  \item Team member names and contributions
\end{itemize}

\subsection*{9.3 Graph Generation}
Students must implement a script to generate performance graphs. Example structure:

\begin{verbatim}
#!/usr/bin/env python3
# Students must implement their own graph generation
# Should read results.csv and produce graphs
import matplotlib.pyplot as plt
import pandas as pd
# Students implement their own plotting logic
\end{verbatim}

Listing 14: Plot Generation Script Structure

\section*{10 Important Notes for Students}
\begin{itemize}
  \item No code implementation is provided - you must write all functions yourself
  \item Only struct definitions and function prototypes are given as guidance
  \item You may discuss algorithms but must write your own code
  \item Plagiarism will result in grade penalty
  \item Each team member must contribute to the implementation
  \item Submit GitHub repository link with all required files
\end{itemize}

\section*{11 Conclusion}
This project requires implementing a complete BZip2-like compression pipeline with configurable parameters, cross-platform support, and comprehensive benchmarking capabilities. Students must write all code themselves based on the provided specifications and prototypes.

\section*{References}
[1] Seward, J. (1996). The bzip2 and libbzip2 official home page.\\[0pt]
[2] Burrows, M., \& Wheeler, D. J. (1994). A block-sorting lossless data compression algorithm.\\[0pt]
[3] Huffman, D. A. (1952). A method for the construction of minimum-redundancy codes.\\[0pt]
[4] Manber, U., \& Myers, G. (1993). Suffix arrays: a new method for on-line string searches.


\end{document}